%! suppress = EnDash
\section{Binary Representation}

\bd[Binary Number]
A \textbf{binary number} is a number expressed in the base-2 numeral system or binary numeral system, a method of
mathematical expression which uses only two symbols: typically $0$ and $1$.
\ed

Binary is useful because it's exactly what you need for representing the values ``true'' and ``false.'' In computers,
an ``on'' state, when electricity is flowing, represents true and the ``off'' state, when no electricity is flowing,
represents false. As in the definition of binary, we can also write 1's and 0's instead of true's and false's. \v

Hence, a single binary value can be used to represent a number. Instead of true and false, we can call these two
states one and zero, which is actually incredibly useful. And if we want to represent larger things, we just need to
add more binary digits. This works exactly the same way as the decimal numbers. With decimal numbers, there are only
ten possible values a single digit can be: 0 through 9. To get numbers larger than 9, we just add more digits to the
front. We can do the same with binary.

\be
For example, let's take the number 263. What does this number actually represent? Well, it means we've got two 100s,
six 10s, and three 1s. If you add those all together, we've got 263. Notice how each column has a different
multiplier; in this case, it's 100, 10, and 1. Each multiplier is 10 times larger than the one to the right. That's
because each column has 10 possible digits to work with, 0 through 9, after which you have to carry one to the next
column. For this reason, it's called base 10 notation, also called decimal since deci means 10.
\ee

Binary works exactly the same way, it's just base 2. That's because there are only two possible digits in binary, 1
and 0. This means that each multiplier has to be two times larger than the column to its right. Instead of 100s, 10s,
and 1s, we now have 4s, 2s, and 1s.

\be
Take for example the binary number 101. This means we have one 4, zero 2s, and one 1. Add those all together, and
we've got the number 5 in base 10.
\ee

But to represent larger numbers, binary needs a lot more digits.

\be
Take this number in binary: 10110111. We can convert it to decimal in the same way. We have:
\bse
1\times128 + 0\times64 + 1\times32 + 1\times16 + 0\times8 + 1\times4 +
1\times2 + 1\times1 = 183
\ese
\ee

Math with binary numbers isn't hard either.

\be
Take for example decimal addition of 183 + 19. First we add 3 + 9 = 12, so we put 2 as the sum and carry 1 to the 10s
column. Now we add 8 + 1 + 1 (that we carried) = 10, so the sum is 0 carry 1. Finally, we add 1 + 1 (that we carried)
= 2, so the total sum is 202. \v

Here's the same sum but in binary. Just as before, we start with the 1s column. Adding 1 + 1 = 2 but there is no
symbol 2, so we use 10 and put 0 as our sum and carry the 1, just like in our decimal example. 1 + 1 + 1 (that we
carried) = 3, or 11 in binary, so we put the sum as 1 and carry 1 again and so on. We end up with this number
11001010, which is the same as the number 202 in base 10.
\ee

Each of these binary digits 1 or 0 is called a bit.

\bd[Bit]
The \textbf{bit}, a contraction of \textbf{bi}nary digi\textbf{t}, is the basic unit of information in computing,
representing a logical state with one of two possible values, most commonly represented as either $1$ or $0$.
\ed

So in these last few examples we were using 8-bit numbers with their lowest value of 0 (i.e.\ $00000000$) and highest
value of 255 (i.e.\ $11111111$). That's 256 different values, or $2^8$. You might have heard of 8-bit computers or
8-bit graphics or audio. These were computers that did most of their operations in chunks of 8 bits, but 256
different values isn't a lot to work with, so it meant things like 8-bit games were limited to just 256 different
colours for their graphics.

\bd[8-bit]
In computer architecture, \textbf{8-bit} integers or other data units are those that are 8 bits wide (1 byte).
\ed

8 bits is such a common size in computing it has a special word: a byte.

\bd[Byte]
The \textbf{byte} is a unit of digital information that consists of eight bits.
\ed

Historically, the byte was the number of bits used to encode a single character of text in a computer and for this
reason it is the smallest addressable unit of memory in many computer architectures. In binary, a kilobyte has
$2^{10}$ bytes, or 1024. \v

There are also 32-bit and 64-bit computer architectures. What this mean is that they operate in chunks of 32 or 64
bits. The largest number you can represent with 32 bits is just under 4.3 billion, which is 32 1s in binary.

\bd[32-bit]
In computer architecture, \textbf{32-bit} integers or other data units are those that are 32 bits wide (4 bytes).
\ed

\bd[64-bit]
In computer architecture, \textbf{64-bit} integers or other data units are those that are 64 bits wide (8 bytes).
\ed

Of course, not everything is a positive number, so we need a way to represent positive and negative numbers. Most
computers use the first bit for the sign, 1 for negative and 0 for positive numbers, and then use the remaining 31
bits for the number itself. That gives us a range of roughly plus or minus two billion. While this is a pretty big
range of numbers, it's not enough for many tasks. This is why 64-bit numbers are useful. The largest value a 64-bit
number can represent is around 9.2 quintillion. Most importantly, computers must label locations in their memory
known as ``addresses'' in order to store and retrieve values. As computer memory has grown to gigabytes and terabytes
– that's trillions of bytes – it was necessary to have 64-bit memory addresses as well. More on that later. \v

In addition to negative and positive numbers, computers must deal with numbers that are not whole numbers, like 12.7
and 3.14. These are called floating-point numbers because the decimal point can float around in the middle of a
number. Several methods have been developed to represent floating-point numbers, the most common of which is the IEEE
754 standard. In essence, this standard stores decimal values sort of like scientific notation. For example, 625.9
can be written as $0.6259 * 10^3$. There are two important numbers here. The 0.6259 is called the ``significand'' and
3 is the ``exponent''. In a 32-bit floating-point number, the first bit is used for the sign of the number, positive
or negative. The next eight bits are used to store the exponent, and the remaining 23 bits are used to store the
significand. \v

Computers use numbers to represent also letters. The most straightforward approach might be to simply number the
letters of the alphabet, A being 1, B being 2, C being 3, and so on. Francis Bacon, the famous English writer, used
5-bit sequences to encode all 26 letters of the English alphabet to send secret messages back in the 1600s. 5 bits
can store 32 possible values, so that's enough for the 26 possible letters but not enough for punctuation, digits,
and upper and lowercase letters. Enter ASCII, the American Standard Code for Information Interchange.

\bd[ASCII]
\textbf{ASCII} abbreviated from American Standard Code for Information Interchange, is a character encoding standard
for electronic communication and is one of the IEEE milestones. ASCII codes represent text in computers,
telecommunications equipment, and other devices.
\ed

Invented in 1963, ASCII was a 7-bit code, enough to store 128 different values. With this expanded range, it could
encode capital letters, lowercase letters, digits 0 through 9, and symbols like the @ sign and punctuation marks.

\be
For example, a lowercase A is represented by the number 97, while a capital A is 65, a colon is 58, and a close
parenthesis is 41.
\ee

ASCII even had a selection of special command codes, such as a newline character to tell a computer where to wrap a
line to the next row. In older computers, the line of text would literally continue off the edge of the screen if you
didn't include a newline character. \v

Because ASCII was such an early standard, it became widely used, and critically allowed different computers built by
different companies to exchange data. This ability to universally exchange information is called ``interoperability''.
However, it did have a major limitation: it was really only designed for English. Fortunately, there are 8 bits in
a byte note 7, and it soon became popular to use codes 128-255, previously unused, for national characters. In the
US, those extra numbers were largely used to encode additional symbols like mathematical notation, graphical
elements, and common accentuated characters. \v

On the other hand, while the Latin characters were used universally, Russian computers used the extra codes to encode
Cyrillic characters and Greek computers used Greek letters and so on. And national character codes worked pretty well
for most countries. The problem was if you opened an email written in Latvian on a Turkish computer, the result was
completely incomprehensible, and things totally broke with the rise of computing in Asia, as languages like Chinese
and Japanese have thousands of characters. There was no way to encode all those characters in 8 bits. In response,
each country invented multibyte encoding schemes, all of which were mutually incompatible. The Japanese were so
familiar with this encoding problem that they even had a special name for it, mojibake, which means ``scrambled
text''. \v

And so it was born ``Unicode'', one format to rule them all.

\bd[Unicode]
\textbf{Unicode}, formally the Unicode Standard, is an information technology standard for the consistent encoding,
representation, and handling of text expressed in most of the world's writing systems. The standard, which is
maintained by the Unicode Consortium, defines 144697 characters covering 159 modern and historic scripts, as well as
symbols, emoji, and non-visual control and formatting codes.
\ed

Devised in 1992 to finally do away with all the different international schemes, it replaced them with one
universal encoding schemes. The most common version of Unicode uses 16 bits with space for over a million codes,
enough for every single character from every language ever used – more than 120000 of them in over 100 types of
script, plus space for mathematical symbols and even graphical characters like emoji. \v

In the same way that ASCII defines a scheme for encoding letters as binary numbers, other file formats like MP3s or
GIFs use binary numbers to encode sounds or colours of a pixel in our photos, movies, and music. Most importantly,
under the hood it all comes down to long sequences of bits. Text messages, this YouTube video, every webpage on the
internet, and even your computer's operating system are nothing but long sequences of 1s and 0s.

\subsection{Logic Gates}

Another reason computers use binary is that an entire branch of mathematics already existed that dealt exclusively
with true and false values. And it had figured out all the necessary rules and operations for manipulating them.
It's called ``Boolean algebra''\footnote{We have introduced a lot of the concepts for Boolean algebra in the first
chapter of the ``Mathematics`` part of the notes however back then we focused on the abstract mathematical part. For
the sake of completeness, we will repeat some of the most basic points again but focusing on the practical aspect.}.

\bd[Boolean Algebra]
\textbf{Boolean algebra}  is the branch of algebra in which the values of the variables are the truth values true and
false, usually denoted 1 and 0, respectively.
\ed

George Boole, from which Boolean Algebra later got its name, was a self-taught English mathematician in the 1800s. He
was interested in representing logical statements that went ``under, over, and beyond'' Aristotle's approach to logic,
which was, unsurprisingly, grounded in philosophy. Boole's approach allowed truth to be systematically and formally
proven, through logic equations which he introduced in his first book ``The Mathematical Analysis of Logic'' in 1847.
In Boolean algebra, the values of variables are true and false, and the operations are logical.

\bd[Logical Operation]
A \textbf{logical operation} is a special symbol or word that connects two or more phrases of information. It is most
often used to test whether a certain relationship between the phrases is true or false.
\ed

\bd[Logic Gate]
A \textbf{logic gate}, or simply \textbf{gate}, is an idealized model of computation or physical electronic device
implementing a logical operation performed on one or more binary inputs that produces a single binary output.
\ed

Logic gates are devices that perform one or all the Boolean logic operations AND, NAND, NOR, NOT, OR, XNOR, and XOR\@.
All types of logic gate, except NOT, accept two binary digits as input, and produce one binary digit as output. NOT
gates accept only one input digit. All digital systems can be constructed by only three basic logic gates. These gates
are the AND gate, the OR gate, and the NOT gate.

\bd[NOT Gate]
A \textbf{NOT gate} (or \textbf{inverter}) is a logic gate which implements logical negation. In mathematical logic
it is equivalent to the logical negation operator (¬).
\ed

A NOT takes a single boolean value, either true or false, and negates it. It flips true to false, and false to true.
We can write out a little logic table that shows the original value under ``Input'', and the outcome after applying the
operation under ``Output''.

\fig{not_gate}{0.5}

\bd[AND Gate]
The \textbf{AND gate} is a basic digital logic gate that implements logical conjunction. In mathematical logic it is
equivalent to the logical conjunction operator (∧).
\ed

AND gate behaves according to the truth table below. A true output (1) results only if all the inputs to the AND gate
are true (1). If none or not all inputs to the AND gate are true, false output results. The function can be extended
to any number of inputs.

\fig{and2}{0.65}

\vspace{-15pt}

\bd[OR Gate]
The \textbf{OR gate} is a basic digital logic gate that implements logical disjunction. In mathematical logic it is
equivalent to the logical disjunction operator (∨).
\ed

OR gate behaves according to the truth table below. A true output (1) results if one or both the inputs to the gate
are true (1). If neither input is true, a false output (0) results. In another sense, the function of OR effectively
finds the maximum between two binary digits, just as the complementary AND function finds the minimum.

\fig{or2}{0.6}

\section[Basic Terminology] {Basic Terminology\footnote{What follows in this section is a brief introduction to some of the most basic terminology used in computer science.
We will not go into much detail, as this is not the purpose of these notes, but we will try to give a brief
introduction to the most important concepts. For a more practical approach, we suggest the reader to consult the
``technical'' part of the notes.}}

\bd[Programming]
\textbf{Programming} is the process of creating a set of instructions that tell a computer how to perform a task.
\ed

Programming is done through the use of ``programming languages''.

\bd[Programming Language]
A \textbf{programming language} is a formal language comprising a set of instructions that produce various kinds of
output, used to implement algorithms.
\ed

The binary representation, is the ``language'' computer processors natively speak, and it is usually called ``machine
code''.

\bd[Machine Code]
\textbf{Machine code} is any low-level programming language, consisting of machine language instructions, which is
used to control a computer's central processing unit.
\ed

In these notes, we will not focus on machine code at all, but we will instead focus on (human-readable) code.
Programming languages are easy and efficient for programmers because they are closer to natural languages than the
machine code. They consist of instructions for computers. These instructions are collected in files usually refered
as ``source code''.

\bd[Source Code]
\textbf{Source code} is any collection of code, with or without comments, written using a human-readable programming
language, usually as plain text.
\ed

The source code of a program is specially designed to facilitate the work of computer programmers, who specify the
actions to be performed by a computer mostly by writing source code. The source code is often transformed by an
assembler or compiler into binary machine code that can be executed by the computer. The machine code might then be
stored for execution at a later time. Alternatively, source code may be interpreted and thus immediately executed. \v

Just like spoken languages, programming languages have statements.

\bd[Statement]
A \textbf{statement} is a syntactic unit of an imperative programming language that expresses some action to be carried
out.
\ed

\be
``A equals 5'' is a programming language statement. In this case, the statement says a variable named $A$ has the number
$5$ stored in it. This is called an assignment statement because we're assigning a value to a variable. \v

To express more complex things, we need a series of statements. For example ``A is 5, B is ten, C equals A plus B''.
This program tells the computer to set variable $A$ equal to $5$, variable $B$ to 10, and finally to add $A$ and $B$
together, and put that result, which is 15, into variable $C$. In general a computer doesn't care about the name of
the variables, as long as variables are uniquely named. But it's probably best practice to name them things that make
sense in case someone else is trying to understand your code.
\ee

A program written in a programming language is formed by a sequence of one or more statements that follow a set of rules
called ``syntax''.

\bd[Syntax]
\textbf{Syntax} is the set of rules that govern the structure and composition of statements in a language.
\ed

A statement may have internal components called ``expressions''.

\bd[Expression]
An \textbf{expression} is a syntactic entity in a programming language that may be evaluated to determine its value.
\ed

An expressions is a combination of one or more constants, operators, and variables, that the programming language
interprets and computes to return another value. Let's see all of these concepts in more detail.

\bd[Constant]
A \textbf{constant} is a value that should not be altered by the program during normal execution, i.e.\ the value is
constant.
\ed

When associated with an identifier, a constant is said to be ``named'', although the terms ``constant'' and ``named
constant'' are often used interchangeably. Constants are useful for both programmers and compilers: For programmers
they are a form of self-documenting code and allow reasoning about correctness, while for compilers they allow
compile-time and run-time checks that verify that constancy assumptions are not violated, and allow or simplify some
compiler optimizations.

\bd[Operator]
An \textbf{operator} is a symbol that tells the compiler to perform specific mathematical or logical manipulations.
\ed

Common simple examples include arithmetic like addition, comparison and greater than, and logical operations like AND\@.
More involved examples include assignments field access in a record or object, and the scope resolution operator.
Languages usually define a set of built-in operators, and in some cases allow users to add new meanings to existing
operators or even define new operators. \v

Constants, functions, and operators are contrasted with a variable, which is an identifier with a value that can be
changed during normal execution, i.e.\ the value is variable.

\bd[Variable]
A \textbf{variable} is an abstract storage location paired with an associated symbolic name, which contains some
known or unknown quantity of information referred to as a value.
\ed

In simple terms, a variable is a container for a particular set of bits or type of data. A variable can eventually be
associated with or identified by a memory address. The variable name is the usual way to reference the stored value,
in addition to referring to the variable itself, depending on the context. This separation of name and content allows
the name to be used independently of the exact information it represents. The identifier in computer source code can
be bound to a value during run time, and the value of the variable may thus change during the course of program
execution.

\bd[Routine / Subroutine / Subprogram/ Function / Method / Procedure]
A \textbf{routine}, or \textbf{subroutine}, ror \textbf{subprogram}, or \textbf{function}, or \textbf{method}, or
\textbf{procedure} is a portion of code within a larger program, which performs a specific task and is relatively
independent of the remaining code.
\ed

Functions may be defined within programs, or separately in libraries that can be used by many programs. Software
consists of thousands of smaller functions, each responsible for different features. In modern programming, it's
uncommon to see functions longer than around 100 lines of code, because by then, there's probably something that
should be pulled out and made into its own function. Modularizing programs into functions not only allows a
single programmer to write an entire app, but also allows teams of people to work efficiently on even bigger programs.
Different programmers can work on different functions, and if everyone makes sure their code works correctly,
then when everything is put together, the whole program should work too! Modern programming languages come with huge
bundles of pre-written functions, called ``libraries''.

\bd[Library]
A \textbf{library} is a collection of resources used by computer programs, often for software development.
\ed

Libraries may include configuration data, documentation, help data, message templates, pre-written code and
subroutines, classes, values or type specifications. Libraries are written by expert coders, made efficient and
rigorously tested,and then given to everyone. There are libraries for almost everything, including networking,
graphics, and sound -- topics we'll discuss in later chapters. For now let's go back to the fundamentals. \v

A program starts at the first statement and runs down one at a time until it hits the end. One can alter this
behaviour by making use of the so called ``control flow statements''.

\bd[Control Flow Statement]
A \textbf{control flow statement} is a statement that results in a choice being made as to which of two or more paths
to follow.
\ed

There are several control flow statement types.

\bd[Conditional Expressions]
\textbf{Conditional expressions} are features of a programming language which perform different computations or
actions depending on whether a programmer-specified boolean condition evaluates to true or false.
\ed

The most used conditional expression is the ``If-Then-Else'' statements which are executed once, a conditional path
is chosen, and the program moves on.

\fig{if}{0.38}

To repeat some statements many times, we need to create a conditional loop.

\bd[Conditional Loops]
\textbf{Conditional loops} are a way for computer programs to repeat one or more various steps depending on
conditions set either by the programmer initially or real-time by the actual program.
\ed

One way is a ``while statement'', also called a ``while loop''. A ``while loop'' checks condition for truthfulness
before executing any of the code in the loop. If condition is initially false, the code inside the loop will never be
executed. \v

There's also the common ``for loop''. Instead of being a condition-controlled loop that can repeat forever until the
condition is false, a ``for loop'' is count-controlled; it repeats a specific number of times. \v

As we mentioned in the previous sections, breaking big programs into smaller functions allows many people to work
simultaneously. They don't have to worry about the whole thing, just the function they're working on. So,  if you're
tasked with writing a sort algorithm, you only need to make sure it sorts properly and efficiently. However, even
packing up code into functions isn't enough. A better solution is to package functions into hierarchies, pulling
related code together into ``objects''.

\bd[Object]
An \textbf{object} is the unit of code that is eventually derived from the process. It can be a variable, a data
structure, a function, or a method.
\ed

\be
For example, car's software might have several functions related to cruise control, like setting speed, nudging speed
up or down, and stopping cruise control altogether. Since they're all related, we can wrap them up into a unified
cruise control object. \v

But, we don't have to stop there. Cruise control is just one part of the engine's software. There might also be sets
of functions that control spark plug ignition, fuel pumps, and the radiator. So we can create a ``parent'' Engine
object that contains all of these ``children'' objects. \v

In addition to children objects, the engine itself might have its own functions. You want to be able to stop and
start it, for example. It'll also have its own variables, like how many miles the car has traveled. In general,
objects can contain other objects, functions, and variables. And of course, the engine is just one part of a Car
object. There's also the transmission, wheels, doors, windows, and so on.
\ee

Programming languages often use something equivalent to the syntax described in the example. The idea of packing up
functional units into nested objects is called Object Oriented Programming.

\bd[Object Oriented Programming (OOP)]
\textbf{Object-oriented programming} (\textbf{OOP}) is a programming paradigm based on the concept of ``objects'',
which can contain data and code.
\ed

OOP is very similar to what we've been doing all long: hide complexity by encapsulating low-level details in
higher-order components. Before we packed up things like transistor circuits into higher-level boolean gates. Now,
we're doing the same thing with software. Yet again, it's a way to move up a new level of abstraction. \v

OOP encapsulates data in the form of fields, often known as ``attributes'' or ``properties'', and code in the form of
procedures, often known as ``methods''.

\bd[Attribute / Property]
An \textbf{attribute} or \textbf{property} is a specification that defines a property of an object.
\ed

For clarity, attributes should more correctly be considered metadata. An attribute is frequently and generally a
property of a property. However, in actual usage, the term attribute can and is often treated as equivalent to a
property depending on the technology being discussed. An attribute of an object usually consists of a name and a value.

\bd[Method]
A \textbf{method} is a function associated with an object.
\ed

Breaking up a big program, like a car's software, into functional units is perfect for teams. One team might be
responsible for the cruise control system, and a single programmer on that team tackles a handful of functions. This
is similar to how big, physical things are built, like skyscrapers. You'll have electricians running wires, plumbers
fitting pipes, welders welding, painters painting, and hundreds of other people teeming all over the hull. They work
together on different parts simultaneously, leveraging their different skills, until one day, you've got a whole
working building. \v

A method cain be either public or private.

\bd[Public Method]
\textbf{Public methods} are methods that are accessible both inside and outside the scope of your class. Any instance
of that class will have access to public methods and can invoke them.
\ed

\bd[Private Method]
\textbf{Private methods} are those methods which can't be accessed in other class except the class in which they are
declared. We can perform the functionality only within the class in which they are declared.
\ed

\be
Returning to our cruise control example: its code is going to have to make use of functions in other parts of the
engine's software to, you know, keep the car at constant speed. That code isn't part of the cruise control team's
responsibility. It's another team's code. Because the cruise control team didn't write that, they're going to need
good documentation about what each function in the code does, and a well-defined Application Programming Interface,
or API for short.
\ee

\bd[Documentation]
\textbf{Documentation} is any communicable material that is used to describe, explain or instruct regarding some
attributes of an object, system or procedure, such as its parts, assembly, installation, maintenance and use.
\ed

Documentation can be provided on paper, online, or on digital or analog media, such as audio tape or CDs. Examples
are user guides, white papers, online help, and quick-reference guides. Paper or hard-copy documentation has become
less common. Documentation is often distributed via websites, software products, and other online applications. \v

Documentation promotes code reuse. So, instead of having programmers constantly write the same things over and over,
they can track down someone else's code that does what they need. Then, thanks to documentation, they can put it to
work in their program, without ever having to read through the code.

\bd[Application Programming Interface (API)]
An \textbf{Application Programming Interface} (\textbf{API}) is a connection between computers or between computer
programs. It is a type of software interface, offering a service to other pieces of software.
\ed

\bd[API Specification]
An \textbf{API specification} is a document or standard that describes how to build or use an API. A computer system
that meets this standard is said to implement or expose an API\@.
\ed

The term API may refer either to the specification or to the implementation. \v

An API connects computers or pieces of software to each other. It is not intended to be used directly by a person
(the end user) other than a computer programmer who is incorporating it into software. An API is often made up of
different parts which act as tools or services that are available to the programmer. A program or a programmer that
uses one of these parts is said to call that portion of the API. The calls that make up the API are also known as
subroutines, methods, requests, or endpoints. An API specification defines these calls, meaning that it explains how
to use or implement them. \v

One purpose of APIs is to hide the internal details of how a system works, exposing only those parts a programmer
will find useful and keeping them consistent even if the internal details later change. An API may be custom-built
for a particular pair of systems, or it may be a shared standard allowing interoperability among many systems.

\be
For example, in the Ignition Control object, there might be functions to set the RPM of the engine, check the spark
plug voltage, as well as fire the individual spark plugs. Being able to set the motor's RPM is really useful. The
cruise control team is going to need to call that function, but they don't know much about how the ignition system
works. It's not a good idea to let them call functions that fire the individual spark plugs, or the engine might
explode! Maybe. The API allows the right people access to the right functions and data.
\ee

OOP languages allow the right people to have access to the right functions and data by letting you specify whether
methods are public or private. This ability to hide complexity, and selectively reveal it,  is the essence of OOP,
and it's a powerful and popular way to tackle building large and complex programs. Pretty much every piece of
software on your computer, or game running on your console was built using an OOP language, like C++, C\#, or
Objective-C. Other popular OOP languages you may have heard of are Python and Java. \v

It's important to remember that code, before being compiled, is just text. As I mentioned earlier, you could write
code in Notepad or any old word processor. Some people do. But generally, today's software developers use
special-purpose applications for writing programs, ones that integrate many useful tools for writing, organizing,
compiling and testing code. Because they've put everything you need in one place, they're called Integrated Development
Environments, or IDEs for short.

\bd[Integrated Development Environment (IDE)]
An \textbf{Integrated Development Environment} (\textbf{IDE}) is a software application that provides comprehensive
facilities to computer programmers for software development.
\ed

An IDE normally consists of at least a source code editor for writing code (often with useful features like automatic
color-coding to improve readability, syntax error checking, and many more), build automation tools and a debugger.
Some IDEs contain the necessary compiler, interpreter, or both; others, such as do not. The boundary between an IDE and
other parts of the broader software development environment is not well-defined; sometimes a version control system or
various tools to simplify the construction of a graphical user interface (GUI) are integrated. Many modern IDEs also
have a class browser, an object browser, and a class hierarchy diagram for use in OOP software development. \v

In addition to coding and debugging, another important part of a programmer's job is documenting their code. This can
be done in standalone files called ``README'' which tell other programmers to read that help file before diving in.

\bd[README]
A \textbf{README} file contains information about the other files in a directory or archive of computer software. A
form of documentation, it is usually a simple plain text file. The file's name is generally written in uppercase.
\ed

It can also happen right in the code itself with comments. These are specially-marked statements that the program
knows to ignore when the code is compiled. They exist only to help programmers figure out what's what in the source
code. Good documentation helps programmers when they revisit code they haven't seen for a while, but it's also crucial
for programmers who are totally new to it.

% TODO: Finish up this section
\section{Algorithms \& Data Structures}

\bd[Data]
\textbf{Data} are the quantities, characters, or symbols on which operations are performed by a computer, which may
be stored and transmitted in the form of electrical signals and recorded on magnetic, optical, or mechanical
recording media.
\ed

If data are arranged in a systematic way then they get a structure and become meaningful. These meaningful (or
processed) data are called ``information''.

\bd[Information]
\textbf{Information} is processed, organized and structured data.
\ed

Information provides context for data and enables decision-making process.

\be
For example, a single customer's sale at a restaurant is data, this becomes information when the business is able to
identify the most popular or least popular dish.
\ee

A data structure is a systematic way to organize data so that they can be used efficiently.

\bd[Data Structure]
A \textbf{data structure} is a data organization, management,and storage format that enables efficient access and
modification.
\ed

More precisely, a data structure is a collection of data values, the relationships among them, and the functions or
operations that can be applied to the data, i.e.\ it is an algebraic structure about data. Data structures are
essential ingredients in creating fast and powerful algorithms, and they make code cleaner and easier to understand. \v

\bd[Abstract Data Type (ADT)]
An \textbf{abstract data type} (\textbf{ADT}) is a mathematical model for data types which is defined by its behavior
(semantics) from the point of view of a user of the data and more specifically in terms of possible values, possible
operations on data of this type, and the behavior of these operations.
\ed

ADTs contrast with data structures, which are concrete representations of data, and are the point of view of an
implementer, not a user. \v

Formally, an ADT may be defined as a ``class of objects whose logical behavior is defined by a set of values and a
set of operations'', which is analogous to an algebraic structure in mathematics. In practice, many common data types
are not ADTs, as the abstraction is not perfect. \v

\textbf{Data structures are used to implements ADTs.} There are multiple ways (i.e.\ data structures) to implement an
ADT. In simple words, ADTs tell us what is to be done and data structures tell us how to do it. In reality, different
implementations of ADTs (a.k.a.\ different data structures) are compared for time and space efficiency. The one best
suited depends on the requirements of the project. \v

Data structures are divided into categories based on their characteristics. One basic division is between ``linear''
and ``non-linear'' data structures.

\bd[Linear Data Structures]
\textbf{Linear data structures} have data elements arranged in linear (sequential) order and each member element is
connected to its previous and next element, except from the first and the last.
\ed

\be
Examples of linear data structures are lists, queues, stacks and arrays.
\ee

\bd[Non-Linear Data Structures]
\textbf{Non-linear data structures} are those data structures in which data items are not arranged in a sequence.
\ed

\be
Examples of non-linear data structures are Trees and Graphs.
\ee

Another important division is between ``static'' and ``dynamic'' data structures.

\bd[Static Data Structures]
A \textbf{static data structure} is a data structure where the data in memory are fixed in size and allocated during
compile time.
\ed

In a static data structure the maximum size is needed to be known in advance, as memory cannot be reallocated at a
later point. Such data structures are easy to implement as computer memory is also sequential. Static data structures
have the advantage of fast access, but they are slower in insertion and deletion.

\be
An example of static data structure is an array.
\ee

\bd[Dynamic Data Structures]
A \textbf{dynamic data structure} is a data structure where the data in memory are allocated in real time.
\ed

A dynamic data structure refers to a collection of data in memory that has the flexibility to grow or shrink in size,
enabling a programmer to control exactly how much memory is utilised. Dynamic data structures have the advantage of
fast insertion and deletion, but they are much slower in access.

\be
An example of dynamic data structure is a linked-list.
\ee

Efficiency of data structures is always measured in terms of time and (memory) space. An ideal data structure would
be the one that takes the least possible time for all its operations and consumes the least memory space. \v

A data type is an attribute of data, as defined by the domain of values data can take, the programming language used,
or the operations that can be performed on the data.

\bd[Data Type]
A \textbf{data type} (or simply \textbf{type}) is an attribute of data which tells the compiler or interpreter how
the programmer intends to use the data. A data type constrains the values that a data might take and the operations
that can be done on the data.
\ed

In other words a data type provides a set of values from which an expression may take its values. \v

Most programming languages support basic data types called ``primitive data types''.

\bd[Primitive Data Type]
\textbf{Primitive data types} are types that are built-in or basic to a language implementation.
\ed

\be
Classic basic primitive types that are generally supported more or less directly by computer hardware and any
programming language are:
\bit
\item \textbf{Character}: is a unit of information that roughly corresponds to a symbol, such as in an alphabet, in the
written form of a natural language.
\item \textbf{Integer}: represents some range of mathematical integers that are commonly represented in a computer as
a group of bits. The size of the grouping varies so the set of integer sizes available varies between different types
of computers.
\item \textbf{Floating-point}: is an arithmetic unit using formulaic representation of real numbers as an approximation
to support a trade-off between range and precision.
\item \textbf{Fixed-point}: refers to a method of representing fractional numbers by storing a fixed number of digits of
their fractional part.
\item \textbf{Boolean}: is a data type that has one of two possible values (usually denoted ``true'' and ``false'')
which is intended to represent the two truth values of logic and Boolean algebra.
\item \textbf{Reference}: is a value that enables a program to indirectly access a particular data, such as a variable's
value or a record, in the computer's memory or in some other storage device.
\eit
\ee

Opposite to primitive data types, we also have the so called ``composite data types``.

\bd[Composite Data Type]
A \textbf{composite data type} is any data type which can be constructed in a program using the programming
language's primitive data types or other composite data types.
\ed